\section*{Bab \ref{ch:b3:forms} --- Bentuk Diferensial dan
Teorema Stokes}

\begin{solution}{exo:b3:forms:1}
Dengan menguraikan dan membunuh faktor berulangnya:
\[
\omega\wedge\eta = (x\,\dd y - z\,\dd x)\wedge(\dd x +
y\,\dd z)
= -x\,\dd x\wedge\dd y + xy\,\dd y\wedge\dd z +
yz\,\dd z\wedge\dd x
\]
(dengan memakai $\dd y\wedge\dd x = -\dd x\wedge\dd y$ dan $-zy\,\dd
x\wedge\dd z = yz\,\dd z\wedge\dd x$). Berikutnya $\dd\omega =
\dd x\wedge\dd y - \dd z\wedge\dd x = \dd x\wedge\dd y +
\dd x\wedge\dd z$ dan $\dd\eta = \dd y\wedge\dd z$. Akhirnya
\[
\dd(\omega\wedge\eta) = y\,\dd x\wedge\dd y\wedge\dd z +
z\,\dd y\wedge\dd z\wedge\dd x = (y + z)\,\dd x\wedge\dd
y\wedge\dd z
\]
(sebab suku pertama $\omega\wedge\eta$ menyumbang $\dd(-x)
\wedge\dd x\wedge\dd y = 0$; sedangkan permutasi siklik tiga
faktor bersifat genap). Pemeriksaan Leibniznya: $\dd\omega\wedge\eta =
(\dd x\wedge\dd y + \dd x\wedge\dd z)\wedge(\dd x + y\,\dd
z) = y\,\dd x\wedge\dd y\wedge\dd z$, dan
$(-1)^1\omega\wedge\dd\eta = -(x\,\dd y - z\,\dd
x)\wedge\dd y\wedge\dd z = z\,\dd x\wedge\dd y\wedge\dd z$;
sehingga jumlahnya cocok.
\end{solution}

\begin{solution}{exo:b3:forms:2}
Bentuk $\dd f = \sum_i\partial_if\,\dd x_i$ berkoefisien
$\nabla f$. Untuk bentuk kerjanya,
\[
\dd\omega_F = (\partial_yF_3 - \partial_zF_2)\,\dd y\wedge\dd
z + (\partial_zF_1 - \partial_xF_3)\,\dd z\wedge\dd x +
(\partial_xF_2 - \partial_yF_1)\,\dd x\wedge\dd y =
\sigma_{\operatorname{curl}F},
\]
yakni bentuk fluks curl-nya (kumpulkanlah keenam suku
$\sum_i\dd F_i\wedge\dd x_i$). Untuk bentuk fluksnya,
$\dd\sigma_F = (\partial_xF_1 + \partial_yF_2 +
\partial_zF_3)\,\dd x\wedge\dd y\wedge\dd z$: yakni
divergensinya. Lalu $\dd^2f = 0$ berbunyi
$\sigma_{\operatorname{curl}\nabla f} = 0$, yakni
$\operatorname{curl}\nabla f = 0$, sedangkan $\dd^2\omega_F = 0$
berbunyi $(\operatorname{div}\operatorname{curl}F)\,\dd
x\wedge\dd y\wedge\dd z = 0$: jadi kedua identitas vektornya adalah
satu identitas, $\dd^2 = 0$, yang dibaca pada dua derajat.
\end{solution}

\begin{solution}{exo:b3:forms:3}
(a) Ketertutupannya adalah kesetangkupan turunan parsial silangnya:
$\partial_y(2xy + z^2) = 2x = \partial_x(x^2)$,
$\partial_z(2xy + z^2) = 2z = \partial_x(2xz)$,
$\partial_z(x^2) = 0 = \partial_y(2xz)$. Sedangkan daerah asalnya $\R^3$
berbentuk bintang: jadi \omterm{def:b3:forms:closedexact}{eksak} (\cref{thm:b3:forms:poincare}), dengan
primitif $f = x^2y + xz^2$ (periksalah $\dd f$).
(b) $\dfrac{x\,\dd x + y\,\dd y}{x^2 + y^2} =
\tfrac12\,\dd\log(x^2 + y^2)$: jadi \omterm{def:b3:forms:closedexact}{eksak} pada seluruh
$\R^2\setminus\{0\}$ (sehingga \omterm{def:b3:forms:closedexact}{tertutup}) --- yakni sepupu radial
bentuk sudutnya yang tak berbahaya.
(c) Pada $\{x > 0\}$ bentuknya adalah $\omega_\theta$, yang \omterm{def:b3:forms:closedexact}{tertutup}
(\cref{exo:b3:forms:4}); sedangkan setengah bidangnya cembung, sehingga ia
\omterm{def:b3:forms:closedexact}{eksak} di sana, dan memang $f = \arctan(y/x)$ memenuhi $\dd f
= \frac{-y\,\dd x + x\,\dd y}{x^2 + y^2}$. Jadi \omterm{def:b3:forms:closedexact}{eksak} pada
setengah bidangnya, tak \omterm{def:b3:forms:closedexact}{eksak} pada bidang tertusuknya: sebab
halangannya tinggal di lubangnya, bukan di rumusnya.
\end{solution}

\begin{solution}{exo:b3:forms:4}
Untuk ketertutupannya: dengan $\rho^2 = x^2 + y^2$,
\[
\partial_x\Bigl(\frac{x}{\rho^2}\Bigr) = \frac{\rho^2 -
2x^2}{\rho^4} = \frac{y^2 - x^2}{\rho^4} =
\partial_y\Bigl(\frac{-y}{\rho^2}\Bigr),
\]
sehingga $\dd\omega_\theta = \bigl(\partial_x(x/\rho^2) -
\partial_y(-y/\rho^2)\bigr)\,\dd x\wedge\dd y = 0$. Lalu pada
$\gamma(t) = (r\cos2\pi t, r\sin2\pi t)$:
\[
\gamma^*\omega_\theta = \frac{2\pi r^2(\cos^22\pi t +
\sin^22\pi t)}{r^2}\,\dd t = 2\pi\,\dd t,
\qquad\text{sehingga}\qquad
\int_\gamma\omega_\theta = 2\pi .
\] Seandainya $\omega_\theta$
\omterm{def:b3:forms:closedexact}{eksak}, integral ini akan lenyap
(\cref{prop:b3:forms:exactloop}): jadi ia tak \omterm{def:b3:forms:closedexact}{eksak}. Seandainya
$\R^2\setminus\{0\}$ berbentuk bintang terhadap suatu $p$,
lema Poincaré (yang digeser ke $p$) akan menjadikan setiap
\omterm{def:b3:forms:closedexact}{bentuk tertutup} \omterm{def:b3:forms:closedexact}{eksak} --- yakni kontradiksi. Secara langsung: bagi sembarang $p
\neq 0$, ruas dari $p$ ke titik $-p$ milik
daerah asalnya melewati $0$: sehingga sifat berbentuk bintangnya gagal pada setiap
titiknya.
\end{solution}

\begin{solution}{exo:b3:forms:5}
Untuk vektor $v_1, \dots, v_{n-1}$, menguraikan determinannya
sepanjang kolom pertamanya memberikan
\[
\det(\nu, v_1, \dots, v_{n-1}) =
\sum_{i=1}^n(-1)^{i-1}\nu_i\,
\det\bigl(\text{baris} \neq i \text{ dari } (v_1, \dots,
v_{n-1})\bigr)
= \sigma_\nu(v_1, \dots, v_{n-1}),
\]
sebab $(\dd x_1\wedge\dots\wedge\widehat{\dd
x_i}\wedge\dots\wedge\dd x_n)(v_1, \dots, v_{n-1})$ tepat
merupakan minor terhapus ke-$i$
(\cref{def:b3:forms:wedge}). Lalu menerapkannya pada $v_j =
\partial_j\gamma$: $(\gamma^*\sigma_\nu)(e_1, \dots,
e_{n-1}) = \det A$ dengan $A = (\nu\ \partial_1\gamma\ \cdots\
\partial_{n-1}\gamma)$. Sekarang ${}^t\!AA$ diagonal per blok:
$\langle\nu, \nu\rangle = 1$ dan $\langle\nu,
\partial_j\gamma\rangle = 0$ (sebab $\nu$ normal, sedangkan
$\partial_j\gamma$ singgung), sehingga $(\det A)^2 =
\det({}^t\!AA) = \det G$; dan $\det A > 0$ bagi parametrisasi
langsung (sebab itulah arti orientasi-$\nu$-nya):
$\gamma^*\sigma_\nu = \sqrt{\det G}\,\dd
u_1\wedge\dots\wedge\dd u_{n-1}$, yakni unsur luas Gramnya. Untuk
$S^2$, $\nu(x) = x$ dan parametrisasi bolanya
$\gamma(\theta, \varphi) = (\sin\varphi\cos\theta,
\sin\varphi\sin\theta, \cos\varphi)$ pada
$(0,2\pi)\times(0,\pi)$ (yang langsung; ia melewatkan satu meridian, yakni
himpunan yang tak berluas): $\det G = \sin^2\varphi$, sehingga
$\int_{S^2}\sigma_\nu =
\int_0^{2\pi}\!\!\int_0^\pi\sin\varphi\,\dd\varphi\,
\dd\theta = 4\pi$.
\end{solution}

\begin{solution}{exo:b3:forms:6}
Bentuk $\omega$ yang diberikan adalah $\sigma_\nu$ bagi $\nu(x) = x$ pada
$S^2$, sehingga (a) adalah perhitungan yang baru saja dikerjakan:
$\int_{S^2}\omega = 4\pi$. (b) $\dd\omega = 3\,\dd
x\wedge\dd y\wedge\dd z$, dan Stokes pada bola satuannya memberikan
$\int_{S^2}\omega = 3\operatorname{vol}(B^3)$: sehingga
$\operatorname{vol}(B^3) = \frac{4\pi}3$. Secara umum,
bentuk $\sigma = \sum_i(-1)^{i-1}x_i\,\dd
x_1\wedge\dots\wedge\widehat{\dd x_i}\wedge\dots\wedge\dd
x_n$ terbatas pada $S^{n-1}$ menjadi bentuk luasnya (dengan $\nu = x$ pada
\cref{exo:b3:forms:5}), $\dd\sigma = n\,\dd
x_1\wedge\dots\wedge\dd x_n$, dan Stokes menghasilkan
$\operatorname{area}(S^{n-1}) = n\operatorname{vol}(B^n)$
--- yang selaras dengan rumus \omterm{ex:b3:lebesgue:gamma}{fungsi Gamma}
\cref{thm:b3:product:ballvolume}.
\end{solution}

\begin{solution}{exo:b3:forms:7}
(a) Stokes yang diterapkan pada bentuk fluks $\sigma_F$ pada
daerah \omterm{def:b3:topology:compact}{kompak} $\Omega$:
$\int_\Omega\operatorname{div}F\,\dd x =
\int_{\partial\Omega}\sigma_F$ (\cref{exo:b3:forms:2} bagi
sisi \omterm{def:b3:topology:interior}{interiornya}). Kenalilah integran batasnya: bagi
$v_1, v_2$ yang singgung pada suatu titik $\partial\Omega$,
uraian kolom pertama \cref{exo:b3:forms:5} memberikan
$\sigma_F(v_1, v_2) = \det(F, v_1, v_2)$; lalu dengan menulis $F =
\langle F, \nu\rangle\nu + T$ dengan $T$ yang singgung, kolom
$T$ merupakan kombinasi bidang $v_1, v_2$, sehingga $\det(T,
v_1, v_2) = 0$ dan
$\sigma_F\vert_{\partial\Omega} = \langle F,
\nu\rangle\,\sigma_\nu = \langle F, \nu\rangle\,\dd S$: yakni
teorema divergensinya, dengan $\nu$ sebagai normal keluarnya
(sebab normal-keluar-dulu tepat merupakan \omterm{def:b3:forms:orientation}{orientasi} terinduksinya).
(b) Stokes yang diterapkan pada $\omega_F$ pada
permukaan-dengan-batas $S$: $\dd\omega_F =
\sigma_{\operatorname{curl}F}$ terbatas menjadi
$\langle\operatorname{curl}F, \nu\rangle\,\dd S$ lewat pengenalan
yang sama, sedangkan pada kurva batasnya
$\gamma^*\omega_F = \langle F\circ\gamma,
\gamma'\rangle\,\dd t$, yakni $\oint\langle F,
\tau\rangle\,\dd\ell$. Pada setengah bola atas dengan $\nu$ yang
keluar (radial): di titik khatulistiwa $p = (1,0,0)$
vektor keluar-di-dalam-permukaannya adalah $-e_3$; lalu melengkapinya
menjadi kerangka positif menunjukkan bahwa khatulistiwanya ditelusuri
berlawanan arah jarum jam dilihat dari atas ($+e_3$): yakni aturan tangan
kanannya, dengan konvensi yang sama pada kedua ruas identitasnya.
\end{solution}

\begin{solution}{exo:b3:forms:8}
Terapkanlah teorema divergensi (\cref{exo:b3:forms:7}(a),
yang buktinya bebas dimensi) pada $F = u\nabla v$:
$\operatorname{div}(u\nabla v) = u\,\Delta v + \langle\nabla
u, \nabla v\rangle$ dan $\langle F, \nu\rangle =
u\,\partial_\nu v$: yakni identitas pertamanya. Lalu menukar $u, v$ lalu
mengurangkannya meniadakan suku setangkupnya: yakni yang kedua. Jika
$\Delta u = 0$ pada $\Omega$ dan $u = 0$ pada $\partial\Omega$:
maka identitas pertamanya dengan $v = u$ memberikan
$\int_\Omega\norm{\nabla u}^2 = 0$, sehingga $\nabla u \equiv 0$
dan $u$ konstan pada setiap komponennya; sedangkan tutupan setiap
komponennya menemui $\partial\Omega$ (menurut keterbatasannya), yang $u = 0$ di sana:
jadi $u \equiv 0$. Lalu dua \omterm{prop:b3:conformal:harmonicholo}{fungsi harmonik} dengan nilai \omterm{def:b3:forms:boundary}{batas}
yang sama berselisih $u$ semacam itu: jadi keduanya berimpit --- yakni ketunggalan
bagi masalah Dirichletnya, yang melengkapi teori keberadaan
pada cakramnya di \cref{ch:b3:conformal}.
\end{solution}

\begin{solution}{exo:b3:forms:9}
Menurut \cref{thm:b3:forms:pullback}(c),
$\varphi^*(f\,\dd y_1\wedge\dots\wedge\dd y_n) =
(f\circ\varphi)\,\det(D\varphi)\,\dd
x_1\wedge\dots\wedge\dd x_n$, sehingga identitas \omterm{def:b3:forms:pullback}{tarikan baliknya} berbunyi
\[
\int_U(f\circ\varphi)\,\det(D\varphi)\,\dd x = \int_Vf\,\dd
y .
\]
Karena $\det D\varphi > 0$ di mana-mana, $\det D\varphi =
\abs{\det D\varphi}$, dan inilah persis rumus
penggantian variabelnya
(\cref{thm:b3:product:changeofvar}) bagi integran \omterm{def:b3:topology:continuity}{kontinu}
yang bertumpuan \omterm{def:b3:topology:compact}{kompak}: jadi masing-masing pernyataannya adalah
yang lain. Adapun nilai mutlaknya masuk ke dalam \emph{hipotesisnya}:
yakni \omterm{def:b3:forms:orientation}{orientasinya}. Bagi $\varphi$ yang membalik \omterm{def:b3:forms:orientation}{orientasi}, identitas
bentuknya memperoleh sebuah tanda minus global (sebab bentuk merasakan
\omterm{def:b3:forms:orientation}{orientasi}), sedangkan rumus \omterm{def:b3:measure:measure}{ukurannya} tetap memakai $\abs{\det}$
(sebab \omterm{def:b3:measure:measure}{ukuran} tidak): jadi dua pembukuan atas satu Jacobi.
\end{solution}

\begin{solution}{exo:b3:forms:10}
Cangkang $A = \{r_1 \leq \norm x \leq r_2\}$ merupakan
submanifold-$3$ \omterm{def:b3:topology:compact}{kompak} dengan \omterm{def:b3:forms:boundary}{batas} $S_{r_2}\cup S_{r_1}$; sedangkan
\omterm{def:b3:forms:orientation}{orientasi} terinduksinya adalah \omterm{def:b3:forms:orientation}{orientasi} sfera yang biasa pada
$S_{r_2}$ (keluar dari $A$ = menjauhi $0$) dan yang
\emph{berlawanan} pada $S_{r_1}$ (keluar dari $A$ = menuju
$0$). Lalu Stokes dengan $\dd\omega = 0$:
\[
0 = \int_A\dd\omega = \int_{S_{r_2}}\omega -
\int_{S_{r_1}}\omega .
\]
Untuk bentuk sudut ruangnya: dengan $\rho = \norm x$ dan $\sigma =
x\,\dd y\wedge\dd z + y\,\dd z\wedge\dd x + z\,\dd
x\wedge\dd y$, berlaku $\omega = \rho^{-3}\sigma$ dan
\[
\dd\omega = \sum_i\partial_i\bigl(x_i\rho^{-3}\bigr)\,
\dd x\wedge\dd y\wedge\dd z
= \bigl(3\rho^{-3} - 3\rho^{-5}\cdot\rho^2\bigr)\dd
x\wedge\dd y\wedge\dd z = 0 .
\]
Pada $S_r$, bagi $v_1, v_2$ yang singgung: $\omega(v_1, v_2) =
r^{-3}\det(x, v_1, v_2) = r^{-3}\,r\det(\nu, v_1, v_2) =
r^{-2}\,\dd S(v_1, v_2)$, sehingga $\int_{S_r}\omega =
r^{-2}\cdot4\pi r^2 = 4\pi$: yakni konstan terhadap $r$, seperti yang diramalkan
deformasinya, dan tak nol --- sehingga $\omega$ \omterm{def:b3:forms:closedexact}{tertutup} tetapi tak
\omterm{def:b3:forms:closedexact}{eksak} pada $\R^3\setminus\{0\}$ (sebab \omterm{def:b3:forms:closedexact}{bentuk eksak} \omterm{def:b3:lebesgue:l1}{terintegralkan} menjadi
$0$ atas $S_r$ yang tanpa \omterm{def:b3:forms:boundary}{batas} menurut Stokes). Inilah hukum
Gauss: bahwa fluks medan sebuah muatan satuan lewat sembarang
sfera yang melingkupinya adalah $4\pi$, berapa pun jari-jarinya.
\end{solution}

\begin{solution}{exo:b3:forms:11}
Tulislah $\omega = c\,\omega_\theta + \dd f$
(\cref{exo:b3:forms:12}) dengan $c =
\frac1{2\pi}\int_{S^1}\omega$. Maka
\[
\int_\gamma\omega = c\int_\gamma\omega_\theta +
\int_\gamma\dd f
= c\cdot2\pi\operatorname{Ind}_\gamma(0) + 0
= n\int_{S^1}\omega,
\]
menurut \cref{prop:b3:forms:exactloop} dan definisi
indeksnya. Satu bilangan bulat $n$ mengendalikan setiap periode pada
bidang tertusuknya: sebab bentuk-$1$ \omterm{def:b3:forms:closedexact}{tertutup} tak dapat membedakan dua
gelung yang berbilangan lilitan sama.
\end{solution}

\begin{solution}{exo:b3:forms:12}
Ambillah $\alpha = \omega - c\,\omega_\theta$: yang \omterm{def:b3:forms:closedexact}{tertutup}, dan
$\int_{S^1}\alpha = 0$ menurut pilihan $c$-nya
(sebab $\int_{S^1}\omega_\theta = 2\pi$). Tariklah balik lewat pemetaan
kutub $\Phi(\rho, \theta) = (\rho\cos\theta, \rho\sin\theta)$,
yakni difeomorfisma lokal surjektif $(0, \infty)\times\R \to
U$: maka $\Phi^*\alpha$ \omterm{def:b3:forms:closedexact}{tertutup}
(\cref{thm:b3:forms:pullback}(b)) pada \omterm{def:b3:topology:topology}{himpunan terbuka} \emph{cembung}
$(0,\infty)\times\R$, sehingga \omterm{def:b3:forms:closedexact}{eksak}
(\cref{thm:b3:forms:poincare}): $\Phi^*\alpha = \dd g$.
Untuk $\rho$ yang tetap: $g(\rho, \theta + 2\pi) - g(\rho, \theta)
= \int_\theta^{\theta + 2\pi}\partial_\theta g\,\dd s$ adalah
integral $\alpha$ mengelilingi lingkaran berjari-jari
$\rho$, yang sama dengan $\int_{S^1}\alpha = 0$ (sebab anulus
antara kedua lingkarannya merupakan permukaan \omterm{def:b3:topology:compact}{kompak} dengan \omterm{def:b3:forms:boundary}{batas};
lalu Stokes seperti pada \cref{exo:b3:forms:10}, satu dimensi lebih rendah).
Jadi $g$ berperiode $2\pi$ terhadap $\theta$ dan turun ke
fungsi $f$ yang terdefinisi baik pada $U$ dengan $f\circ\Phi = g$;
lalu $f$ mulus (sebab $\Phi$ difeomorfisma lokal) dan
$\Phi^*(\dd f) = \dd g = \Phi^*\alpha$ memaksa $\dd f =
\alpha$. Karena itu $\omega = c\,\omega_\theta + \dd f$.
Untuk ketunggalannya: mengintegralkannya atas $S^1$ menetapkan $c$, sebab \omterm{def:b3:forms:closedexact}{bentuk
eksak} berperiode nol; sedangkan $f$ tunggal sampai penambahan
konstanta. Jadi pemetaan $[\omega] \mapsto
\frac1{2\pi}\int_{S^1}\omega$ merupakan isomorfisma
linear $H^1(\R^2\setminus\{0\}) \to \R$, dan kelas
$\omega_\theta$ membangkitkannya: sehingga lubangnya tepat
berdimensi satu, secara kohomologis.
\end{solution}

\begin{solution}{pb:b3:forms:1}
\textbf{1.} $\dd\sigma = \sum_i(-1)^{i-1}\,\dd x_i\wedge\dd
x_1\wedge\dots\wedge\widehat{\dd x_i}\wedge\dots\wedge\dd
x_n$; lalu membawa $\dd x_i$ melintasi $i-1$ faktor sebelumnya
memakan $(-1)^{i-1}$, yang meniadakan prafaktornya: sehingga masing-masing dari
$n$ sukunya sama dengan $\dd x_1\wedge\dots\wedge\dd x_n$, jadi
$\dd\sigma = n\,\dd x_1\wedge\dots\wedge\dd x_n$. Lalu Stokes pada
bolanya: $\int_S\sigma = \int_{\bar B}\dd\sigma =
n\operatorname{vol}(\bar B) > 0$.

\textbf{2.} $n = 2$: $\sigma = x\,\dd y - y\,\dd x$; lalu pada
$\gamma(t) = (\cos t, \sin t)$, $\gamma^*\sigma = (\cos^2t +
\sin^2t)\,\dd t = \dd t$, yakni bentuk panjang busurnya:
$\int_S\sigma = 2\pi = 2\operatorname{vol}(\bar B^2)$.
$n = 3$: $\sigma\vert_S$ adalah bentuk luasnya
(\cref{exo:b3:forms:5} dengan $\nu(x) = x$):
$\int_S\sigma = 4\pi = 3\cdot\tfrac{4\pi}3$. Keduanya cocok dengan
pertanyaan 1.

\textbf{3.} Dengan menurunkan $\norm\varphi^2 = 1$: $2\langle
D\varphi(x)h, \varphi(x)\rangle = 0$ bagi setiap $h$, sehingga
$\operatorname{im}D\varphi(x) \subseteq \varphi(x)^\perp$, yakni sebuah
hiperbidang: jadi $\operatorname{rk}D\varphi(x) \leq n - 1$. Karena
$\dd\sigma$ merupakan bentuk-$n$ (pertanyaan 1), maka titik demi titik
$\varphi^*(\dd\sigma)_x =
(D\varphi(x))^*(\dd\sigma)_{\varphi(x)} = 0$ menurut
\cref{prop:b3:forms:linearpullback}(2): sebab pemetaan ke dalam
sferanya tak punya ruang untuk menarik balik sebuah volume.

\textbf{4.} Cacatnya ada pada ``sehingga'' yang kedua: sebab pada
himpunan berdimensi $(n-1)$ yaitu $S$, setiap bentuk-$(n-1)$ \omterm{def:b3:forms:closedexact}{tertutup} secara
trivial (karena tak ada bentuk-$n$ tak nol pada
manifold-$(n-1)$), tetapi \emph{\omterm{def:b3:forms:closedexact}{tertutup} tak berarti \omterm{def:b3:forms:closedexact}{eksak}},
sedangkan $\int_S\dd\eta = 0$ menuntut primitif $\eta$ yang sungguh-sungguh
terdefinisi pada $S$. Padahal pembatasan $\sigma$ justru tak
\omterm{def:b3:forms:closedexact}{eksak} --- sebab integralnya $n\operatorname{vol}(\bar B) \neq
0$ --- dan ketakeksakan inilah yang menggerakkan seluruh soalnya.

\textbf{5.} Kita akan mengintegralkan $r^*\sigma$ atas sferanya
lalu menghitungnya dengan dua cara. Karena $r$ menetapkan $S$ titik demi titik,
integralnya sama dengan $\int_S\sigma = n\operatorname{vol}(\bar B)
\neq 0$. Sedangkan karena $r$ terdefinisi pada bolanya, Stokes mengubah
integral yang sama itu menjadi $\int_{\bar B}\dd(r^*\sigma) =
\int_{\bar B}r^*(\dd\sigma)$; dan karena $r$ bernilai
di sferanya, pertanyaan 3 membuat integran itu lenyap. Jadi satu
bilangan, dua nilai: sehingga retraksinya tak mungkin ada.

\textbf{6.} Kedua integralnya dihitung lewat parametrisasi
langsung $\gamma$ atas keping-keping $S$
(\cref{def:b3:forms:integral}); dan karena $r\circ\gamma =
\gamma$ (sebab parametrisasinya mendarat di $S$, yang $r$ menjadi
identitas di sana), maka $\gamma^*(r^*\sigma) = (r\circ\gamma)^*\sigma =
\gamma^*\sigma$ (\cref{thm:b3:forms:pullback}(a)): sehingga integran
lokalnya berimpit, dan \omterm{lem:b3:forms:partition}{partisi kesatuan} mana pun memberikan
$\int_Sr^*\sigma = \int_S\sigma$.

\textbf{7.} Bentuk $r^*\sigma$ adalah bentuk-$(n-1)$ mulus pada
\omterm{def:b3:topology:topology}{persekitaran} $\bar B$ yang \omterm{def:b3:topology:compact}{kompak} dan berorientasi, yang
batasnya dengan \omterm{def:b3:forms:orientation}{orientasi} terinduksinya adalah $S$: sehingga Stokes
(\cref{thm:b3:forms:stokes}) memberikan persis
$\int_Sr^*\sigma = \int_{\bar B}\dd(r^*\sigma)$.

\textbf{8.} Berlaku $\dd(r^*\sigma) = r^*(\dd\sigma)$
(\cref{thm:b3:forms:pullback}(b)), yang lenyap menurut
pertanyaan 3 yang diterapkan pada $\varphi = r$. Lalu merangkai pertanyaan
6--8:
\[
0 < n\operatorname{vol}(\bar B) = \int_S\sigma =
\int_Sr^*\sigma = \int_{\bar B}\dd(r^*\sigma) = 0 :
\]
yakni mustahil. \emph{Tak ada retraksi mulus $\bar B^n$
ke $S^{n-1}$} ($n \geq 2$).

\textbf{9.} Sebuah $r\colon\intcc{-1}1\to\{-1,1\}$ yang \omterm{def:b3:topology:continuity}{kontinu}
dengan $r(\pm1) = \pm1$ akan memetakan himpunan \omterm{def:b3:topology:connected}{terhubung} ke
$\{-1, 1\}$ yang tak \omterm{def:b3:topology:connected}{terhubung}, dan itu mustahil: sebab peta \omterm{def:b3:topology:continuity}{kontinu}
himpunan \omterm{def:b3:topology:connected}{terhubung} bersifat \omterm{def:b3:topology:connected}{terhubung} --- setara dengan itu,
teorema nilai antaranya akan memaksa $r$ mengambil nilai
$0$. Pernyataan yang sama pada setiap dimensi persis adalah Bagian
II; sedangkan keterhubungannya adalah bayangan berdimensi $1$ bagi
halangan kohomologis $\int_S\sigma \neq 0$.

\textbf{10.} Fungsi $x \mapsto \norm{g(x) - x}$ \omterm{def:b3:topology:continuity}{kontinu} dan
di mana-mana $> 0$ pada $\bar B$ yang \omterm{def:b3:topology:compact}{kompak}: sehingga minimumnya
$\delta$ tercapai, jadi $> 0$.

\textbf{11.} Persamaan $\norm{x + tu}^2 = 1$ berbunyi $t^2 + 2t\langle x,
u\rangle + \norm x^2 - 1 = 0$, yang akarnya adalah
\[
t_\pm = -\langle x, u\rangle \pm \sqrt{\langle x,
u\rangle^2 + 1 - \norm x^2} .
\]
Hasil kalinya $\norm x^2 - 1 \leq 0$: sehingga akarnya mengapit
$0$ (atau salah satunya lenyap), jadi parameter sinarnya --- yakni
akar tak negatifnya --- adalah $t(x) = t_+$.

\textbf{12.} Seandainya isi akarnya lenyap di $x \in \bar B$:
maka karena kedua sukunya tak negatif, $\norm x = 1$ dan
$\langle x, u(x)\rangle = 0$, yakni $\langle x, x -
g(x)\rangle = 0$, yakni $\langle x, g(x)\rangle = 1$. Lalu menurut
Cauchy--Schwarz, $1 = \langle x, g(x)\rangle \leq
\norm x\,\norm{g(x)} \leq 1$: yakni kesamaan sepanjang rantainya, yang
memaksa $g(x)$ segaris dengan $x$, bernorma $1$, dan searah:
jadi $g(x) = x$ --- yang terkecualikan. Jadi isi akarnya \omterm{def:b3:topology:continuity}{kontinu} dan
$> 0$ pada $\bar B$, sehingga terbatas di bawah oleh suatu $c > 0$
di sana dan pada sebuah \omterm{def:b3:topology:topology}{persekitaran} (menurut \omterm{def:b3:topology:continuity}{kekontinuan} seragamnya). Pada
\omterm{def:b3:topology:topology}{persekitaran} itu, $u$ mulus (sebab $\norm{x - g(x)} \geq
\delta/2$, dengan mengecilkannya bila perlu), isi akarnya tetap $\geq
c/2$, dan akar kuadratnya mulus pada $\intoo0\infty$: sehingga $r$
mulus di dekat $\bar B$.

\textbf{13.} Berlaku $\norm{r(x)} = 1$ menurut konstruksi $t(x)$-nya:
jadi $r(\bar B) \subseteq S$. Untuk $x \in S$: $\langle x,
u(x)\rangle = \frac{1 - \langle x, g(x)\rangle}{\norm{x -
g(x)}} \geq 0$ (menurut Cauchy--Schwarz sekali lagi), dan $\norm x =
1$ menyusutkan isi akarnya menjadi $\langle x, u\rangle^2$, yang
akar kuadratnya adalah $\langle x, u\rangle$ itu sendiri (sebab ia $\geq
0$): sehingga $t(x) = 0$ dan $r(x) = x$. Jadi $r$ merupakan retraksi
mulus dari bolanya ke sferanya --- yang bertentangan dengan
Bagian II. \emph{Setiap pemetaan diri mulus $\bar B$ mempunyai
titik tetap.}

\textbf{14.} Berlaku $\varepsilon > 0$ persis seperti pada pertanyaan 10.
Sedangkan polinomialnya membentuk subaljabar $\mathcal C(\bar B,
\R)$ yang memuat konstantanya dan memisahkan titik (sebab $x
\mapsto x_i$ memisahkan), sehingga Stone--Weierstrass
(\cref{thm:b3:complete:stoneweierstrass}) menghampiri setiap
koordinatnya: pilihlah polinomial $p_i$ dengan $\sup_{\bar
B}\abs{p_i - f_i} < \frac{\varepsilon}{2\sqrt n}$; maka
pemetaan vektor $p = (p_1, \dots, p_n)$ memenuhi
$\sup_{\bar B}\norm{p - f} \leq
\bigl(\sum_i\sup_{\bar B}\abs{p_i - f_i}^2\bigr)^{1/2} <
\varepsilon/2$.

\textbf{15.} Pada $\bar B$: $\norm p \leq \norm f +
\frac\varepsilon2 \leq 1 + \frac\varepsilon2$, sehingga
$\norm{g} = \frac{\norm p}{1 + \varepsilon/2} \leq 1$:
jadi $g(\bar B) \subseteq \bar B$. Lebih jauh $\norm{g - p} =
\frac{\varepsilon/2}{1 + \varepsilon/2}\norm p \leq
\frac\varepsilon2$, sehingga $\norm{g - f} \leq \norm{g - p} +
\norm{p - f} < \varepsilon$ pada $\bar B$.

\textbf{16.} Pemetaan $g$ polinomial, jadi mulus, dan memetakan
$\bar B$ ke dirinya sendiri: sehingga Bagian III menyediakan $x_0 = g(x_0)$.
Lalu $\norm{f(x_0) - x_0} = \norm{f(x_0) - g(x_0)} <
\varepsilon = \min_{\bar B}\norm{f - \operatorname{id}}$:
yakni kontradiksi. \emph{Setiap \omterm{def:b3:topology:continuity}{pemetaan kontinu} $\bar B^n \to
\bar B^n$ mempunyai titik tetap.}

\textbf{17.} \emph{Bola terbuka:} $f(x) = \frac{x + e_1}2$ memetakan
$B$ ke dalam $B$ (sebab $\norm{f(x)} < 1$ sejati) sedangkan satu-satunya titik
tetapnya $e_1$ terletak pada sferanya: jadi kekompakannya hilang.
\emph{Sfera:} pemetaan antipodal $x \mapsto -x$ bebas
titik tetap; sebab $S$ \omterm{def:b3:topology:compact}{kompak} tetapi bertopologi ``keliru''
--- yakni justru bukan-retrak Bagian II.
\emph{Anulus:} rotasi sebesar sudut $\not\equiv 0$ mana pun tak menetapkan
apa pun; sebab lubangnya melindungi rotasinya --- jadi kecembungannya (atau lebih
tepatnya, \omterm{def:b3:topology:topology}{topologi} serupa-bolanya) hilang. Jadi teorema Brouwer
sesungguhnya tentang himpunan \emph{cembung \omterm{def:b3:topology:compact}{kompak}}, seperti yang dimanfaatkan
pertanyaan 18.

\textbf{18.} Untuk $x \in \Delta$: ada $x_j > 0$, sehingga $(Ax)_i
\geq A_{ij}x_j > 0$ bagi setiap $i$; jadi $\norm{Ax}_1 > 0$
dan $T(x) = Ax/\norm{Ax}_1$ terdefinisi baik, \omterm{def:b3:topology:continuity}{kontinu}, serta
mendarat di $\Delta$ (sebab entrinya positif dan berjumlah $1$).
Sedangkan $\Delta$ cembung, \omterm{def:b3:topology:compact}{kompak}, dan berinterior tak kosong pada
hiperbidang afin $\{\sum x_i = 1\} \cong \R^{n-1}$; lalu
pemetaan radial dari barisentrumnya --- sebab setiap sinar dari
barisentrumnya menemui $\partial\Delta$ pada tepat satu titik, menurut
kecembungan dan kekompakannya, sedangkan fungsi tolok
yang bersesuaian bersifat \omterm{def:b3:topology:continuity}{kontinu} --- merupakan \omterm{def:b3:topology:continuity}{homeomorfisma} $\Delta \to
\bar B^{n-1}$. Lalu mengangkut Brouwer lewatnya: $T$ mempunyai
titik tetap $x^*$, yakni $Ax^* = \lambda x^*$ dengan $\lambda
= \norm{Ax^*}_1 > 0$; dan $x^* = Ax^*/\lambda$ berentri
positif sejati menurut perhitungan pembukanya. \emph{Jadi
matriks positif mempunyai vektor eigen positif.}

\textbf{19.} Pertanyaan 18 memberikan $A\pi = \lambda\pi$, $\pi \in
\Delta$, $\pi > 0$. Lalu jumlahkanlah koordinatnya: $\sum_i(A\pi)_i =
\sum_j\pi_j\sum_iA_{ij} = \sum_j\pi_j = 1$ (sebab kolomnya berjumlah
$1$), sedangkan $\sum_i\lambda\pi_i = \lambda$. Jadi $\lambda = 1$
dan $A\pi = \pi$: yakni vektor peluang stasioner --- yaitu
kesetimbangan rantai Markovnya, inti matematis PageRank.

\textbf{20.} Pada $S$: $\norm{F_t(x)}^2 = \norm x^2 +
2t\langle x, v(x)\rangle + t^2\norm{v(x)}^2 = 1 + 0 + t^2$
menurut kesinggungannya dan $\norm v = 1$: jadi $F_t(S) \subseteq
\sqrt{1+t^2}\,S$.

\textbf{21.} Tetapkanlah sebuah atlas berhingga berisi parametrisasi
langsung $\gamma$ beserta \omterm{lem:b3:forms:partition}{partisi kesatuan}, yang tak bergantung pada $t$. Pada sebuah
peta, $F_t\circ\gamma = \gamma + t(v\circ\gamma)$, sehingga setiap
koefisien $(F_t\circ\gamma)^*\sigma$ merupakan jumlah
hasil kali satu faktor $(x_i + tv_i)\circ\gamma$ (yang afin terhadap
$t$) dan sebuah determinan $(n-1)\times(n-1)$ yang entrinya
afin terhadap $t$: yakni polinomial dalam $t$ berderajat $\leq n$ dengan
koefisien yang mulus terhadap variabel petanya. Lalu mengalikannya dengan
fungsi partisinya yang tak bergantung pada $t$ dan mengintegralkannya
suku demi suku: $P(t) = \sum_{k=0}^n c_kt^k$, yakni sebuah polinomial.

\textbf{22.} \emph{Untuk keinjektifannya:} $v$ bersifat Lipschitz pada $S$
(sebab mulus pada \omterm{def:b3:topology:compact}{kompak}), katakanlah berkonstanta $L$; lalu untuk $\abs t <
1/L$, $\norm{F_t(x) - F_t(y)} \geq (1 - \abs
tL)\norm{x - y} > 0$. \emph{Untuk difeomorfismanya:} $G_t =
F_t/\sqrt{1 + t^2}$ memetakan $S$ ke $S$; sedangkan pada peta, Jacobinya
konvergen seragam ke Jacobi $G_0 = \operatorname{id}$ saat
$t \to 0$, sehingga untuk $t$ yang kecil semuanya terbalikkan dan $G_t$ merupakan
difeomorfisma lokal (\cref{thm:b3:submanifolds:ift} pada
peta), injektif, dengan peta yang terbuka (menurut difeomorfisma lokalnya) dan
\omterm{def:b3:topology:compact}{kompak} di $S^{n-1}$ yang \emph{\omterm{def:b3:topology:connected}{terhubung}} ($n \geq 2$):
sehingga petanya $= S$, jadi $G_t$ difeomorfisma $S$.
\emph{Untuk penskalaannya:} di bawah $s_\lambda(x) = \lambda x$, setiap
koefisien $x_i$ memperoleh $\lambda$ dan masing-masing dari $n - 1$
diferensialnya memperoleh $\lambda$: jadi $s_\lambda^*\sigma =
\lambda^n\sigma$. Karena $F_t = s_{\sqrt{1+t^2}}\circ G_t$:
\[
P(t) = \int_SG_t^*\bigl(s_{\sqrt{1+t^2}}^{\,*}\sigma\bigr)
= (1 + t^2)^{n/2}\int_SG_t^*\sigma
= (1 + t^2)^{n/2}\int_S\sigma
\quad\text{untuk nilai kecil } t,
\]
dengan kesamaan terakhirnya menurut \cref{lem:b3:forms:consistency} (sebab $G_t$
difeomorfisma $S$ yang mengawetkan \omterm{def:b3:forms:orientation}{orientasi} untuk $t$ yang
kecil: determinan Jacobi petanya berubah secara \omterm{def:b3:topology:continuity}{kontinu},
tak pernah lenyap, dan positif di $t = 0$).

\textbf{23.} Jika $n$ ganjil dan sebuah medan singgung satuan mulus
ada, maka pertanyaan 21--22 membuat polinomial
$P(t)/\int_S\sigma$ (yang sahih: sebab $\int_S\sigma \neq 0$ menurut
pertanyaan 1) berimpit dengan $(1 + t^2)^{n/2}$ di dekat $0$; sedangkan dua
fungsi mulus yang berimpit di dekat $0$ dengan salah satunya polinomial
memaksa $(1 + t^2)^{n/2}$ menjadi polinomial itu pada seluruh
$\R$. Padahal jika $Q(t)^2 = (1 + t^2)^n$ dengan $Q \in \R[t]$,
maka pemfaktoran tunggal di $\R[t]$
(\cref{ch:b3:rings}) memberikan $t^2 + 1$ yang \omterm{def:b3:rings:divisibility}{tak tereduksi} itu
multiplisitas genap pada $Q^2$ dan multiplisitas ganjil $n$ pada
$(1 + t^2)^n$: yakni mustahil. Jadi tak ada medan singgung satuan mulus
pada $S^{n-1}$ untuk $n$ yang ganjil --- yakni sfera
berdimensi genap. Akhirnya, medan singgung yang sekadar \omterm{def:b3:topology:continuity}{kontinu} dan
tak lenyap di mana pun, sebutlah $w$, akan menghasilkan satu: perluaslah ia ke sebuah
\omterm{def:b3:topology:topology}{persekitaran} lewat $\tilde w(x) = w(x/\norm x)$, muluskanlah
komponen demi komponen (\cref{thm:b3:lp:regularization}) menjadi
$v_0$ yang mulus dengan $\sup_S\norm{v_0 - \tilde w} <
\frac12\min_S\norm w$, proyeksikanlah secara singgung $v_1(x) = v_0(x)
- \langle v_0(x), x\rangle x$ --- sebab pada $S$ ini mengubah $v_0$
paling banyak sebesar komponen normalnya, yang sendirinya paling banyak $\norm{v_0 -
\tilde w}$ karena $\tilde w$ singgung, sehingga $\norm{v_1 - \tilde
w} \leq 2\norm{v_0 - \tilde w} < \min\norm w$ dan $v_1$
tak pernah lenyap pada $S$ --- lalu normalkanlah: sehingga $v =
v_1/\norm{v_1}$ merupakan medan singgung satuan yang mulus. Karena itu pada
setiap sfera berdimensi genap, \emph{setiap medan singgung
\omterm{def:b3:topology:continuity}{kontinu} mempunyai nol}: sehingga setiap angin di Bumi menyisakan satu titik
tenang.

\textbf{24.} Berlaku $\langle v(x), x\rangle = \sum_j(-y_jx_j +
x_jy_j) = 0$: jadi singgung; dan $\norm{v(x)} = \norm x = 1$: jadi satuan.
Kemulusannya jelas (sebab pemetaan linear). Jadi setiap sfera
berdimensi ganjil mengemban medan singgung satuan yang mulus --- yakni perkalian
oleh $\iu$ sepanjang garis kompleksnya --- sehingga
halangan pertanyaan 23 tepat berupa paritas dimensinya. Pada
argumen polinomialnya, untuk $n$ yang genap fungsi $(1 +
t^2)^{n/2}$ memang \emph{adalah} polinomial, dan tak ada kontradiksi
yang muncul: jadi buktinya bukan sekadar gagal berlaku, melainkan
kesimpulannya sungguh salah, seperti yang disaksikan $v$.

\textbf{25.} Andaikan $f$ tak pernah lenyap pada $\bar B$. Maka
$g(x) = -f(x)/\norm{f(x)}$ \omterm{def:b3:topology:continuity}{kontinu} $\bar B \to S
\subseteq \bar B$, dan Brouwer (pertanyaan 16) menyediakan $x^* =
g(x^*)$. Karena $g$ bernilai di $S$, maka $x^* \in S$, dan
\[
\langle f(x^*), x^*\rangle = \bigl\langle f(x^*),\
-\tfrac{f(x^*)}{\norm{f(x^*)}}\bigr\rangle
= -\norm{f(x^*)} < 0,
\]
yang bertentangan dengan hipotesis batasnya. Jadi $f$ mempunyai nol.
Untuk kesurjektifannya: diberikan $y \in \R^n$, terapkanlah yang di atas pada $f(x) =
F(x) - y$ pada bola $\bar B(0, R)$ dengan $R$ yang begitu besar sehingga
$\langle F(x), x\rangle \geq \norm y\,\norm x$ pada sfera
berjari-jari $R$ (menurut koersivitasnya); maka $\langle f(x), x\rangle =
\langle F(x), x\rangle - \langle y, x\rangle \geq 0$ di sana
(menurut Cauchy--Schwarz), sehingga pernyataan yang diskalakan ulang itu memberikan nol
bagi $f$: yakni $F(x) = y$. Jadi setiap medan \omterm{def:b3:topology:continuity}{kontinu} yang koersif bersifat pada
--- yakni bayangan bebas derajat bagi teorema keberadaan
variasionalnya, yang diantarkan \omterm{def:b3:topology:topology}{topologi} murni.

\end{solution}
